\documentclass[11pt]{article}

\usepackage[margin=1in]{geometry}
\usepackage{amsmath, amssymb, amsthm}
\usepackage{bm}
\usepackage{hyperref}
\usepackage{enumitem}
\usepackage{graphicx}
\usepackage{mathtools}
\usepackage{booktabs}
\usepackage{listings}
\usepackage[T1]{fontenc}
\usepackage{lmodern}

\hypersetup{
  colorlinks=true,
  linkcolor=blue,
  urlcolor=blue,
  citecolor=blue
}

\lstset{
  basicstyle=\ttfamily\small,
  breaklines=true,
  frame=single,
  columns=fullflexible,
  keepspaces=true
}

\title{Zeta Genius and Z-MOD:\\
A Multiplicity-Theoretic, Zeta-Guided Optimization Framework\\
for Analytic Multiplicity Probing}
\author{Citizen Gardens \\ The Foundation of Multiplicity}
\date{April 28, 2026}

\begin{document}
\maketitle

\begin{abstract}
This defensive publication describes the Zeta Genius project and the Z-MOD (Zeta-Mathematics Optimizer with Multiplicity Dynamics) framework as developed across multiple design and analysis steps. Z-MOD is a hybrid continuous--discrete optimizer that embeds neural network parameters into the complex plane via a multiplicity-theoretic mapping, couples updates to the Riemann zeta function and its zeros, and incorporates a discrete prime-multiplicity state. The document formalizes the geometric embedding, zeta potentials, zero-data-informed resonance layers, and multiplicity dynamics. It further specifies minimal PyTorch implementations, benchmark protocols on MNIST and CIFAR-10, an ADR-based development roadmap, and analytic probes for distinguishing zeta-specific effects from generic structured oscillations. The intent is to create prior art around this family of constructions, preventing later patenting while enabling open research and extension.
\end{abstract}
\newpage
\tableofcontents

\section{Executive Summary}

This publication records and formalizes a set of interconnected ideas and concrete artifacts developed under the ``Zeta Genius'' and ``Z-MOD'' umbrella.

At a high level, the contributions are:
\begin{itemize}[leftmargin=*]
  \item A \emph{multiplicity-aware} optimizer state space
  \[
    \mathcal{S} = \mathbb{R}^n \times \mathcal{M}, \qquad \mathcal{M} = \{0,1\}^m,
  \]
  where $\theta \in \mathbb{R}^n$ are continuous parameters and $\bm{\pi} \in \mathcal{M}$ encodes discrete prime multiplicity modes.
  \item A \emph{canonical geometric embedding} into the complex plane
  \[
    s(\theta, \bm{\pi}) = \sigma_0 + i\left(a\,\tilde{\phi}(\theta) + \sum_{p\in\mathcal{P}} \pi_p \log p\right),
  \]
  with $\tilde{\phi}$ a clipped parameter-energy functional and $\mathcal{P}$ a small prime set, e.g.\ $\{2,3,5\}$.
  \item A \emph{zeta-guided potential} based on the Riemann zeta function
  \[
    \mathcal{Z}_0(\theta,\bm{\pi}) = \Re \log\big(|\zeta(s(\theta,\bm{\pi}))| + \varepsilon\big),
  \]
  and its exact logarithmic derivative $\zeta'/\zeta$ pulled back through the embedding to define analytic gradients.
  \item A \emph{zero-data-informed resonance layer} built from tabulated imaginary parts of zeta zeros, introducing resonance bands and explicit-formula-inspired oscillations without assuming the Riemann Hypothesis:
  \[
    \Delta \mathcal{Z}_{\mathrm{zeros}}(\theta,\bm{\pi}) = \lambda_{\mathrm{zeros}} \frac{1}{N}\sum_{n=1}^N
      \exp\left(-\frac{\big(\mathrm{Im}\,s(\theta,\bm{\pi}) - \gamma_n\big)^2}{2\sigma^2}\right).
  \]
  \item A \emph{multiplicity dynamics layer} that evolves $\bm{\pi}$ via a symmetric Metropolis--Hastings kernel on $\{0,1\}^m$, coupled to the zeta-informed objective.
  \item A \emph{minimal PyTorch implementation} of ZMODAdam, a research optimizer that wraps Adam with togglable zeta scaling, auxiliary zeta losses, zero-resonance terms, and multiplicity flips, along with MNIST and CIFAR-10 training harnesses.
  \item A \emph{phased ADR (Architecture Decision Record) roadmap and repository scaffold} for systematically developing, validating, and scaling Z-MOD, including specific benchmark matrices and logging formats.
  \item An \emph{analytic multiplicity probe} design: toy Rastrigin experiments and MNIST ablations intended to distinguish three cases:
  \begin{enumerate}
    \item Case A: no effect of zeta-specific structure.
    \item Case B: structured oscillation helps, but zeta zeros are not special.
    \item Case C: zeta-specific advantages that survive random-band and generic-oscillation controls.
  \end{enumerate}
\end{itemize}

The combination of geometric embedding, zeta potentials, prime multiplicity, and zero-data-informed resonance---implemented in a concrete optimizer and wrapped in an ADR-governed roadmap---constitutes the core novelty of this prior art.

\section{Background and Prior Art}

\subsection{Classical optimization and adaptive methods}

Modern deep learning relies heavily on gradient-based optimizers such as SGD, momentum, RMSProp, and Adam, all of which operate on $\theta \in \mathbb{R}^n$ without an explicit number-theoretic or complex-analytic structure. Adam and its variants maintain exponential moving averages of gradients and squared gradients and have become near-default optimizers for many tasks.

There also exist optimizers and algorithms that incorporate notions of sharpness (e.g.\ SAM, ``Sharpness-Aware Minimization'') and noise injections, with empirical gains in generalization. These methods shape the effective loss surface via local perturbation or reweighting, but they typically do not explicitly reference zeta functions or prime distributions.

\subsection{Zeta-based optimization proposals}

Recent work such as ZetA (``Zeta Algorithm'') introduces scalar zeta-based scaling of gradients and loss terms, reporting empirical gains on certain tasks. These designs typically:
\begin{itemize}[leftmargin=*]
  \item Use $\zeta$ as a scalar scaling factor in an otherwise standard optimizer (often Adam).
  \item Combine zeta scaling with additional regularizers (e.g.\ entropy, sharpness penalties).
\end{itemize}
As a result, the role of zeta-specific structure is entangled with other components, making it difficult to isolate whether the arithmetic properties of $\zeta$ contribute anything beyond general structured perturbations.

\subsection{Multiplicity theory and prime encoding}

The Zeta Genius project is grounded in a multiplicity-theoretic viewpoint: mathematical objects and algorithmic states are modeled as patterns of prime-labeled interactions and recurrence across scales. In this view, primes play the role of discrete modes or ``colors'' that modulate continuous dynamics, and complex functions such as $\zeta$ encode spectral information about these multiplicity patterns.

The Z-MOD framework is designed precisely to make this multiplicity structure explicit in the optimizer state:
\[
  \mathcal{S} = \mathbb{R}^n \times \mathcal{M}, \qquad \mathcal{M} = \{0,1\}^m,
\]
with $\mathcal{M}$ representing a finite prime-multiplicity subspace.

\subsection{Novelty of Z-MOD}

Compared to prior art:

\begin{itemize}[leftmargin=*]
  \item Z-MOD embeds optimizer parameters into the complex plane via a \emph{multiplicity-aware} mapping $s(\theta,\bm{\pi})$, not merely through a scalar mapping of training step or loss.
  \item Z-MOD introduces a discrete prime multiset $\bm{\pi}$ and couples it to $\zeta$ and its zeros, yielding a hybrid continuous--discrete process.
  \item Z-MOD defines and implements \emph{zero-data-informed resonance layers} based on tabulated zeros, and systematically compares them against random-band and generic-oscillation controls.
  \item The framework explicitly structures the development through ADRs, benchmark matrices, and analytic probes designed to falsifiably test zeta-specific contributions.
\end{itemize}

To the author's knowledge, no prior optimizer:

\begin{enumerate}
  \item Embeds $(\theta, \bm{\pi})$ into the complex plane in this multiplicity-theoretic manner, nor
  \item Uses tabulated zeros and GUE-like spacing heuristics as auxiliary losses and proposal reweighting within a gradient-based optimizer, nor
  \item Wraps such a design in a systematic ADR-based roadmap with concrete PyTorch artefacts as specified here.
\end{enumerate}

\section{Mathematical Framework}

\subsection{State space and embedding}

Let $n \in \mathbb{N}$ be the number of continuous parameters and let $\mathcal{P} = \{p_1,\dots,p_m\}$ be a fixed finite set of primes, e.g.\ $\{2,3,5\}$. Define the discrete prime-multiplicity space
\[
  \mathcal{M} = \{0,1\}^m,
\]
and the full state space
\[
  \mathcal{S} = \mathbb{R}^n \times \mathcal{M}.
\]

A point in $\mathcal{S}$ is denoted $(\theta,\bm{\pi})$ with $\theta \in \mathbb{R}^n$ and $\bm{\pi} = (\pi_{p_1},\dots,\pi_{p_m}) \in \{0,1\}^m$.

Define a base energy functional
\[
  \phi(\theta) = \frac{1}{n}\|\theta\|^2,
\]
and a clipped version
\[
  \tilde{\phi}(\theta) = \min\big(\phi(\theta), \Phi_{\max}\big),
\]
for some fixed $\Phi_{\max} > 0$.

The canonical complex embedding is
\begin{equation}
  s(\theta,\bm{\pi}) = \sigma_0 + i\left(a\,\tilde{\phi}(\theta) + \sum_{p\in\mathcal{P}} \pi_p \log p\right),
  \label{eq:embedding}
\end{equation}
with parameters $\sigma_0 \in (0,1)$ and $a>0$. The imaginary part,
\[
  t(\theta,\bm{\pi}) = \mathrm{Im}\, s(\theta,\bm{\pi}) = a\,\tilde{\phi}(\theta) + \sum_{p\in\mathcal{P}} \pi_p \log p,
\]
acts as a scalar summary of both continuous and discrete state.

\subsection{Base zeta potential and gradient}

Define the base zeta potential by
\begin{equation}
  \mathcal{Z}_0(\theta,\bm{\pi}) = \Re \log\big(|\zeta(s(\theta,\bm{\pi}))| + \varepsilon\big),
  \label{eq:zeta_potential_base}
\end{equation}
where $\varepsilon>0$ is a small constant ensuring numerical stability.

Let $f(\theta)$ be the task loss (e.g.\ cross-entropy) and define a prime-complexity penalty
\[
  C(\bm{\pi}) = \sum_{p\in\mathcal{P}} \pi_p \log p.
\]

The full objective is
\begin{equation}
  \mathcal{O}(\theta,\bm{\pi}) = f(\theta) + \lambda \mathcal{Z}_0(\theta,\bm{\pi}) + \beta C(\bm{\pi}),
  \label{eq:objective_base}
\end{equation}
with weights $\lambda,\beta \ge 0$.

For $\tilde{\phi}(\theta) < \Phi_{\max}$, the chain rule yields an analytic gradient for the zeta term. Let $s = s(\theta,\bm{\pi})$, and assume $\zeta(s)\neq 0$. Then
\[
  \frac{\partial \mathcal{Z}_0}{\partial \phi}
  = \Re\left(\frac{\zeta'(s)}{\zeta(s)} \cdot i a\right)\cdot \frac{|\zeta(s)|}{|\zeta(s)|+\varepsilon}.
\]
Because $\tilde{\phi}$ is quadratic in $\theta$, we have
\[
  \nabla_\theta \tilde{\phi}(\theta) = \frac{2}{n}\theta,
\]
and thus
\begin{equation}
  \nabla_\theta \mathcal{Z}_0(\theta,\bm{\pi}) =
  \begin{cases}
    \displaystyle
    \Re\Big(\frac{\zeta'(s)}{\zeta(s)} \cdot i a\Big)\,
    \frac{|\zeta(s)|}{|\zeta(s)|+\varepsilon}\,
    \frac{2}{n}\theta, & \phi(\theta) < \Phi_{\max}, \\[8pt]
    0, & \phi(\theta) \ge \Phi_{\max}.
  \end{cases}
  \label{eq:grad_zeta}
\end{equation}

The base gradient of $\mathcal{O}$ with respect to $\theta$ is then
\begin{equation}
  \nabla_\theta \mathcal{O}(\theta,\bm{\pi})
    = \nabla_\theta f(\theta) + \lambda \nabla_\theta \mathcal{Z}_0(\theta,\bm{\pi}).
  \label{eq:grad_objective_base}
\end{equation}

\subsection{Zero-data-informed resonance layer}

To incorporate information about zeta zeros without assuming the Riemann Hypothesis, we introduce a resonance layer based on tabulated zeros.

Let $\{\gamma_n\}_{n=1}^N$ be the imaginary parts of the first $N$ non-trivial zeros of $\zeta(s)$, i.e.\ zeros of the form $\rho_n = \beta_n + i\gamma_n$ with $0<\beta_n<1$ and $0<\gamma_1\le \gamma_2\le\cdots$. These are obtained from numerical tables and are treated as empirical data.

Define a Gaussian-band kernel over the imaginary axis:
\begin{equation}
  R(t) = \frac{1}{N} \sum_{n=1}^N \exp\left(-\frac{(t-\gamma_n)^2}{2\sigma^2}\right),
  \label{eq:kernel_R}
\end{equation}
with $\sigma>0$ controlling band width.

We then define the zero-resonance auxiliary loss
\begin{equation}
  \mathcal{L}_{\mathrm{zero}}(\theta,\bm{\pi}) = \alpha_{\mathrm{zero}}\;\tilde{\phi}(\theta)\; R(t(\theta,\bm{\pi})),
  \label{eq:aux_zero_loss}
\end{equation}
with $\alpha_{\mathrm{zero}}>0$ a small hyperparameter. This term increases when $\theta$ has large energy and $\mathrm{Im}\,s$ lies near a zero ordinate, encouraging the optimizer to explore and respond differently near these resonant bands.

Additionally, to mimic explicit-formula-type oscillations, one may define
\begin{equation}
  E(\theta,\bm{\pi}) = \sum_{n=1}^N w_n \cos\big(\gamma_n \log x(\theta,\bm{\pi})\big),
  \quad x(\theta,\bm{\pi}) = \exp\big(t(\theta,\bm{\pi})\big),
  \label{eq:explicit_style}
\end{equation}
with decaying weights $w_n$ (e.g.\ $w_n = 1/n$) and a separate coefficient $\lambda_{\mathrm{EF}}$.

The total zeta-informed potential becomes
\begin{equation}
  \mathcal{Z}_{\mathrm{inf}}(\theta,\bm{\pi})
    = \mathcal{Z}_0(\theta,\bm{\pi})
      + \Delta \mathcal{Z}_{\mathrm{zeros}}(\theta,\bm{\pi})
      + \lambda_{\mathrm{EF}} E(\theta,\bm{\pi}),
  \label{eq:zeta_informed_total}
\end{equation}
where $\Delta \mathcal{Z}_{\mathrm{zeros}}$ is identified with $\mathcal{L}_{\mathrm{zero}}$ or a closely related formulation.

\subsection{Multiplicity dynamics}

The prime-multiplicity state $\bm{\pi}\in\{0,1\}^m$ evolves via a Markov chain with Metropolis--Hastings acceptance:

\begin{itemize}[leftmargin=*]
  \item Proposal: select index $k\in\{1,\dots,m\}$ uniformly at random and define
  \[
    \bm{\pi}' = \bm{\pi} + e_k \mod 2,
  \]
  where $e_k$ is the unit vector and addition is modulo 2.
  \item Acceptance probability:
  \[
    \alpha\big((\theta,\bm{\pi}) \to (\theta,\bm{\pi}')\big)
      = \min\big(1, \exp(-\beta_{\mathrm{MH}}(\mathcal{O}(\theta,\bm{\pi}') - \mathcal{O}(\theta,\bm{\pi})))\big),
  \]
  where $\beta_{\mathrm{MH}}>0$ is an inverse temperature.
\end{itemize}

The resulting dynamics is a piecewise-deterministic process: $\theta$ evolves via gradient-based updates on $\mathcal{O}$ (or a variant thereof), while $\bm{\pi}$ is updated intermittently via Metropolis proposals.

\subsection{Full objective with auxiliary layers}

Combining the components, the most complete objective in this publication has the form
\begin{equation}
  \mathcal{O}_{\mathrm{full}}(\theta,\bm{\pi})
   = f(\theta) + \lambda \mathcal{Z}_0(\theta,\bm{\pi})
     + \lambda_{\mathrm{zeros}} \tilde{\phi}(\theta) R(t(\theta,\bm{\pi}))
     + \lambda_{\mathrm{EF}} E(\theta,\bm{\pi})
     + \beta C(\bm{\pi}).
  \label{eq:full_objective}
\end{equation}

By toggling each coefficient independently, one can construct ablations that isolate:
\begin{itemize}[leftmargin=*]
  \item Scalar zeta scaling ($\lambda\neq 0$, other terms off),
  \item Auxiliary zeta loss ($\lambda\neq 0$ in combination with selected terms),
  \item Zero-resonance effects ($\lambda_{\mathrm{zeros}}\neq 0$),
  \item Explicit-formula-style oscillations ($\lambda_{\mathrm{EF}}\neq 0$),
  \item Multiplicity dynamics ($\beta>0$ together with Metropolis updates).
\end{itemize}

\section{Algorithmic Design and Code Snippets}

\subsection{ZMODAdam optimizer skeleton}

The core optimizer described in the development thread is a research-oriented variant of Adam, augmented with zeta and multiplicity features. In Python-like pseudocode:

\begin{lstlisting}[language=Python, caption={ZMODAdam skeleton with zero-resonance auxiliary loss}]
class ZMODAdam(torch.optim.Optimizer):
    def __init__(self, params, lr=1e-3,
                 use_zeta_scaling=False,
                 use_aux_zeta_loss=False,
                 use_prime_flips=False,
                 use_zero_resonance=False,
                 aux_zeta_weight=0.0,
                 zero_resonance_weight=0.0,
                 zero_band_sigma=5.0,
                 sigma0=0.6, a=0.1, phi_max=100.0,
                 primes=(2, 3, 5)):
        defaults = dict(lr=lr)
        super().__init__(params, defaults)
        self.use_zeta_scaling = use_zeta_scaling
        self.use_aux_zeta_loss = use_aux_zeta_loss
        self.use_prime_flips = use_prime_flips
        self.use_zero_resonance = use_zero_resonance
        self.aux_zeta_weight = aux_zeta_weight
        self.zero_resonance_weight = zero_resonance_weight
        self.zero_band_sigma = zero_band_sigma
        self.sigma0 = sigma0
        self.a = a
        self.phi_max = phi_max
        self.primes = torch.tensor(primes, dtype=torch.float64)
        # Preload zeros as a 1D tensor: gamma_1,...,gamma_N
        self.zero_gammas = load_riemann_zeros_tensor()
        # Adam state initialization omitted for brevity

    def compute_phi_tensor(self):
        sq_sum = 0.0
        count = 0
        for group in self.param_groups:
            for p in group["params"]:
                if p.grad is None:
                    continue
                sq_sum = sq_sum + (p.data ** 2).sum()
                count += p.data.numel()
        if count == 0:
            return torch.tensor(0.0, device=self.zero_gammas.device)
        phi = sq_sum / float(count)
        return torch.clamp(phi, max=self.phi_max)

    def _prime_shift(self, pi_state):
        # pi_state: tensor of shape [len(primes)] with 0/1 entries
        return torch.dot(pi_state.to(self.primes.dtype), torch.log(self.primes))

    def compute_auxiliary_loss(self, pi_state):
        params = [p for g in self.param_groups for p in g["params"] if p.requires_grad]
        if not params:
            return torch.tensor(0.0, device=self.zero_gammas.device)

        device = params[0].device
        loss = torch.zeros((), device=device)

        # Zero-resonance auxiliary term
        if self.use_zero_resonance and self.zero_resonance_weight != 0.0:
            phi = self.compute_phi_tensor()
            imag_s = self.a * phi + self._prime_shift(pi_state.to(device))
            gammas = self.zero_gammas.to(device)
            diff = imag_s - gammas
            sigma = self.zero_band_sigma
            kernel = torch.exp(-0.5 * (diff / sigma) ** 2)
            R = kernel.mean()
            loss = loss + self.zero_resonance_weight * phi * R

        # Additional auxiliary zeta or oscillation terms can be added here

        return loss

    @torch.no_grad()
    def step(self, closure=None, pi_state=None):
        # closure returns loss; pi_state is the current multiplicity vector
        loss = None
        if closure is not None:
            loss = closure()

        # Optionally add auxiliary loss
        if pi_state is None:
            pi_state = torch.zeros(len(self.primes), dtype=torch.int64)
        aux_loss = self.compute_auxiliary_loss(pi_state)
        if aux_loss.requires_grad:
            aux_loss.backward()

        # Standard Adam update logic here (omitted for brevity)

        return loss
\end{lstlisting}

This skeleton emphasizes the independence of features: zeta scaling, auxiliary losses, and zero resonance are all toggled by explicit flags.

\subsection{MNIST training harness}

A minimal training script for MNIST classification using ZMODAdam might look as follows:

\begin{lstlisting}[language=Python, caption={MNIST training harness with ZMODAdam}]
def train_mnist(model, optimizer, train_loader, device):
    model.train()
    total_loss = 0.0
    correct = 0
    total = 0
    for x, y in train_loader:
        x, y = x.to(device), y.to(device)
        optimizer.zero_grad()
        logits = model(x)
        loss = F.cross_entropy(logits, y)
        loss.backward()
        optimizer.step()
        total_loss += loss.item() * x.size(0)
        preds = logits.argmax(dim=1)
        correct += (preds == y).sum().item()
        total += x.size(0)
    return total_loss / total, correct / total

def eval_mnist(model, data_loader, device):
    model.eval()
    total_loss = 0.0
    correct = 0
    total = 0
    with torch.no_grad():
        for x, y in data_loader:
            x, y = x.to(device), y.to(device)
            logits = model(x)
            loss = F.cross_entropy(logits, y)
            total_loss += loss.item() * x.size(0)
            preds = logits.argmax(dim=1)
            correct += (preds == y).sum().item()
            total += x.size(0)
    return total_loss / total, correct / total

def main():
    # Data, model, optimizer setup
    device = torch.device("cuda" if torch.cuda.is_available() else "cpu")
    model = MNISTMLP().to(device)
    optimizer = ZMODAdam(
        model.parameters(), lr=1e-3,
        use_zeta_scaling=False,
        use_aux_zeta_loss=False,
        use_prime_flips=False,
        use_zero_resonance=True,
        zero_resonance_weight=1e-4,
        zero_band_sigma=5.0
    )
    # Training loop
    for epoch in range(1, 21):
        train_loss, train_acc = train_mnist(model, optimizer, train_loader, device)
        test_loss, test_acc = eval_mnist(model, test_loader, device)
        print(f"Epoch {epoch}: "
              f"train_loss={train_loss:.4f}, train_acc={train_acc:.3f}, "
              f"test_loss={test_loss:.4f}, test_acc={test_acc:.3f})")
\end{lstlisting}

Additional logging (e.g.\ $\mathrm{Im}\,s$, resonance scores, $\bm{\pi}$ occupancy) can be added as described in the roadmap.

\section{Benchmark Protocol and Analytic Multiplicity Probe}

\subsection{Ablation matrix on MNIST}

To empirically isolate the effect of zeta-specific structure, the framework defines a benchmark matrix on MNIST with (at minimum) the following methods:

\begin{center}
\begin{tabular}{ll}
\toprule
Block & Method \\
\midrule
Base & Adam \\
Control & Adam + matched Gaussian gradient noise \\
Control & Adam + periodic scalar modulation \\
Zeta & Adam + scalar zeta scaling \\
Geometry & Adam + auxiliary zeta loss (no multiplicity) \\
Multiplicity & ZMODAdam + prime flips (no zero resonance) \\
Resonance & ZMODAdam + zero-resonance layer (real zeros) \\
Control & ZMODAdam + random-band resonance (fake zeros) \\
\bottomrule
\end{tabular}
\end{center}

Each method is run for $5$ seeds with fixed data pipeline, batch size, learning rate schedule, and number of epochs.

\subsection{Metrics and statistical tests}

Primary metrics:
\begin{itemize}[leftmargin=*]
  \item Test accuracy as a function of epoch.
  \item Training loss as a function of epoch.
  \item Wall-clock time per epoch and per full run.
  \item Approximate ``escape time'': the epoch index at which test accuracy first exceeds a threshold (e.g.\ 98\%).
\end{itemize}

Secondary metrics:
\begin{itemize}[leftmargin=*]
  \item Generalization gap (train vs test loss).
  \item Mean and variance of $\mathrm{Im}\,s$ across epochs.
  \item Mean resonance score $\tilde{\phi} R(t)$.
  \item Prime occupancy statistics (frequency of $\pi_p = 1$ in final states).
\end{itemize}

The primary analytic multiplicity probe on MNIST uses these metrics to distinguish:
\begin{itemize}[leftmargin=*]
  \item Case A: zero-resonance (real zeros) no better than random-band and baseline.
  \item Case B: zero-resonance and random-band both improve over baseline but are indistinguishable from each other.
  \item Case C: zero-resonance improves over both random-band and baseline with statistically significant differences (e.g.\ Mann--Whitney $p<0.01$).
\end{itemize}

\subsection{Toy Rastrigin probe}

In addition to MNIST, a 2D Rastrigin-based probe is defined:

\begin{itemize}[leftmargin=*]
  \item Landscape: $f(\theta)$ is the 2D Rastrigin function.
  \item Optimizers: SGD, SGD+noise, generic sine-oscillation baseline, Z-MOD without resonance, Z-MOD with resonance.
  \item Metrics: number of iterations to reach $f(\theta)< \varepsilon$, $\mathrm{Im}\,s(t)$ trajectories, prime occupancy, and resonance scores.
\end{itemize}

This toy probe can be run quickly on a CPU to obtain early Case A/B/C indications before moving to MNIST and CIFAR-10.

\section{ADR Roadmap and Repository Scaffold}

\subsection{Phased ADR roadmap}

To ensure disciplined development, the project is organized into phases with associated ADRs, including (non-exhaustively):

\begin{itemize}[leftmargin=*]
  \item ADR-000: Project scope and falsifiability.
  \item ADR-001: Metrics and evidence policy.
  \item ADR-010: Canonical embedding specification.
  \item ADR-020: ZMODAdam core design.
  \item ADR-030: Prime multiset dynamics.
  \item ADR-040: Zero-resonance layer (zero-data-informed).
  \item ADR-050: MNIST benchmark matrix.
  \item ADR-051: CIFAR-10 benchmark matrix.
  \item ADR-060: Reproducibility and artifact packaging.
\end{itemize}

Each ADR documents:
\begin{itemize}[leftmargin=*]
  \item The decision and context.
  \item Alternatives considered.
  \item Rationale and consequences.
  \item Validation plan and file impacts.
\end{itemize}

\subsection{Repository structure}

A recommended repository structure is:

\begin{verbatim}
zmod/
├── README.md
├── pyproject.toml
├── requirements.txt
├── data/
│   └── zeros/
│       ├── riemann_zeros_first_100.txt
│       └── riemann_zeros_first_1000.txt
├── configs/
│   ├── mnist_base.yaml
│   ├── mnist_zeta_scaling.yaml
│   ├── mnist_zero_resonance.yaml
│   ├── cifar10_base.yaml
│   └── cifar10_zmod.yaml
├── docs/
│   ├── adrs/
│   └── validation/
├── src/
│   └── zmod/
│       ├── embedding.py
│       ├── multiplicity.py
│       ├── resonance.py
│       ├── logging_utils.py
│       ├── metrics.py
│       ├── models/
│       └── optim/
├── scripts/
│   ├── train_mnist.py
│   ├── train_cifar10.py
│   ├── run_mnist_sweep.py
│   ├── run_cifar10_sweep.py
│   └── plot_benchmarks.py
├── tests/
├── results/
└── reports/
\end{verbatim}

This structure aligns code, configs, data, tests, and documentation around the Z-MOD design, easing reproducibility and extension.

\section{Discussion and Recommendations}

\subsection{Defensive publication scope}

The key claim of this defensive publication is not that Z-MOD is superior to all existing optimizers, but that the following design space is now prior art:

\begin{itemize}[leftmargin=*]
  \item Optimizers that embed $(\theta,\bm{\pi})$ into the complex plane via a multiplicity-weighted mapping.
  \item Use of $\zeta$ and $\zeta'/\zeta$ as potential terms and gradient modifiers in such embeddings.
  \item Use of small prime sets and Boolean multiplicity vectors as auxiliary optimizer state.
  \item Use of tabulated zeta zeros and their empirical statistics as auxiliary losses and proposal biases via resonance bands and explicit-formula-style oscillations.
  \item Systematic ADR-governed development with benchmark protocols designed to test zeta-specificity via ablations.
\end{itemize}

Any future attempt to patent this family of methods should be precluded by the existence of this description and associated code.

\subsection{Future work}

The current design suggests several natural extensions:

\begin{itemize}[leftmargin=*]
  \item Larger $\mathcal{P}$ and richer multiplicity spaces (beyond $\{2,3,5\}$).
  \item Alternative embeddings (layer-wise energies, gradient-based energies, or spectral norms).
  \item Incorporation of higher-level number-theoretic functions (e.g.\ Dirichlet $L$-functions) for specialized tasks.
  \item Application to zeta-algorithm families such as prime prediction, encryption, signal processing, Schrödinger dynamics, and chaos models.
\end{itemize}

Any such extensions, if they reuse the core pattern of $(\theta,\bm{\pi})$ embedding, zeta potentials, and zero-data-informed resonance, should be understood as lying within the conceptual scope of this defensive publication.

\section{Finite-Width Spectral Effects, RH Deviations, and Convergence Limits in Z-MOD}
\label{sec:finite_width_rh_convergence}

This section analyzes how finite-width spectral effects, deviations from the Riemann Hypothesis (RH), and explicit-formula truncation constrain what can be claimed about the convergence of Z-MOD. The key message is that convergence must be justified by analytic properties of the implemented kernels and standard stochastic-approximation arguments, not by Montgomery's pair correlation conjecture or RH themselves.

\subsection{Montgomery Pair Correlation and Z-MOD Dynamics}

Montgomery's pair correlation conjecture (PCC) describes the limiting two-point correlation of normalized ordinates of non-trivial zeros of the Riemann zeta function on the critical line.\footnote{See, e.g., Montgomery's original work and expository sources for the precise statement and its GUE interpretation.} After appropriate normalization to unit mean spacing, the conjectured pair correlation function is
\[
  R_2(u) = 1 - \left(\frac{\sin(\pi u)}{\pi u}\right)^2,
\]
which is identical to the sine-kernel law for eigenvalues of large Gaussian Unitary Ensemble (GUE) matrices.

Z-MOD does not directly operate on the zero set of $\zeta$; instead, it induces a ``spectral'' process via the imaginary part of its embedding
\[
  t_k = \operatorname{Im}s(\theta_k,\bm{\pi}_k),
\]
where $(\theta_k,\bm{\pi}_k)$ is the optimizer state at iteration $k$. The resonance layer uses a finite subset $\{\gamma_n\}_{n=1}^N$ of zero ordinates as parameters in smooth kernels (e.g., Gaussian bands or cosine sums), but the dynamics of $\{t_k\}$ is governed by both these kernels and the underlying task loss $f(\theta)$.

\begin{lemma}[PCC does not imply Z-MOD convergence]
Even under RH and the truth of Montgomery's PCC, no convergence guarantee for Z-MOD follows directly from PCC. The PCC is an asymptotic statistical statement about normalized zero spacings; Z-MOD uses a finite, truncated, and smoothed kernel built from a finite prefix of zero ordinates, and its state evolution is driven by gradient descent on a task-dependent objective rather than by sampling from the zero point process.
\end{lemma}

\begin{proof}
Montgomery's PCC concerns the limit as the height $T\to\infty$ of the pair correlation of zeros in the range $[0,T]$, after an explicit normalization. By contrast, Z-MOD uses a fixed finite set $\{\gamma_n\}_{n=1}^N$ and constructs a smooth kernel such as
\[
  R_N(t) = \frac{1}{N}\sum_{n=1}^N \exp\!\left(-\frac{(t-\gamma_n)^2}{2\sigma^2}\right).
\]
The optimizer evolution is then
\[
  \theta_{k+1} = \theta_k - \eta_k\left(\nabla f(\theta_k) + g(\theta_k,\bm{\pi}_k) + \xi_k\right),
\]
where $g$ encodes zeta and resonance contributions and $\xi_k$ is stochastic noise. The distribution of $\{t_k\}$ is a complicated function of $f$, $\{\gamma_n\}$, the kernel, and the noise; it is not a sample from the zero process. Thus the asymptotic pair correlation statistics of zeros do not, by themselves, control the trajectory of $\{\theta_k\}$ or imply any almost-sure convergence to critical points of an objective. Convergence must instead be derived from local regularity and step-size conditions for the specific vector field used by Z-MOD.
\end{proof}

\subsection{RH Deviations and Lipschitz Stability}

Many explicit-formula constructions naturally assume that zeros lie on the critical line, i.e., that each non-trivial zero has the form $\rho_n = \tfrac12 + i\gamma_n$. To model possible deviations from RH, write
\[
  \rho_n = \tfrac12 + \delta_n + i\gamma_n,
\]
with $\delta_n \in \mathbb{R}$ describing horizontal deviations from the critical line. In a truncated explicit-formula-inspired auxiliary loss of the form
\[
  \mathcal{L}_{\mathrm{VM}}(\theta,\bm{\pi})
  = \lambda_{\mathrm{VM}}
    \sum_{n=1}^N w_n \Re\!\left(\frac{x(\theta,\bm{\pi})^{\rho_n}}{\rho_n}\right),
  \qquad x(\theta,\bm{\pi}) = e^{t(\theta,\bm{\pi})},
\]
these deviations appear as amplitude factors $x^{\delta_n}$.

\begin{lemma}[Lipschitz constant under horizontal deviations]
Suppose $t(\theta,\bm{\pi})$ is restricted to a compact interval $[-T_{\max}, T_{\max}]$, so that $x(\theta,\bm{\pi}) = e^{t(\theta,\bm{\pi})}$ lies in $[e^{-T_{\max}}, e^{T_{\max}}]$. Assume $|\delta_n|\le \Delta$ for all zeros used in $\mathcal{L}_{\mathrm{VM}}$. Then the Lipschitz constant of $\nabla_\theta \mathcal{L}_{\mathrm{VM}}$ on any compact set in parameter space grows at most by a factor of $e^{\Delta T_{\max}}$ relative to the RH-idealized case $\delta_n=0$.
\end{lemma}

\begin{proof}
For each term,
\[
  \Re\!\big(x^{\rho_n}\big) = x^{1/2 + \delta_n} \cos(\gamma_n\log x),
\]
so $|\Re(x^{\rho_n})| \le x^{1/2 + |\delta_n|}$. On the domain $x\in[e^{-T_{\max}}, e^{T_{\max}}]$, we have $x^{|\delta_n|} \le e^{|\delta_n| T_{\max}} \le e^{\Delta T_{\max}}$. Thus, compared to the case $\delta_n=0$, each term’s magnitude is inflated by at most $e^{\Delta T_{\max}}$. The gradient $\nabla_\theta \mathcal{L}_{\mathrm{VM}}$ involves a common factor of
\[
  \sum_{n=1}^N w_n \Re\!\big(x^{\rho_n}\big)
\]
times $\nabla_\theta t(\theta,\bm{\pi})$, which is independent of $\delta_n$. Therefore the operator norm and the Lipschitz constant of $\nabla_\theta \mathcal{L}_{\mathrm{VM}}$ are inflated by at most the same factor $e^{\Delta T_{\max}}$, up to fixed multiplicative constants depending on $\{w_n\}$ and the bounds on $\nabla_\theta t$ on the compact set.
\end{proof}

This lemma illustrates why clipping $\tilde{\phi}(\theta)$ and thereby controlling the range of $t(\theta,\bm{\pi})$ is critical: it bounds the potential impact of RH deviations on the regularity of the auxiliary gradient field.

\subsection{Spectral Leakage in Truncated Resonance Bands}

The resonance layer uses a truncated zero set to build a smooth kernel
\[
  R_N(t) = \frac{1}{N}\sum_{n=1}^N K_\sigma(t - \gamma_n),
  \qquad K_\sigma(u) = \exp\!\left(-\frac{u^2}{2\sigma^2}\right),
\]
where $\{\gamma_n\}_{n=1}^N$ are the imaginary parts of the first $N$ non-trivial zeros and $\sigma>0$ is a bandwidth parameter. Let $\mu$ be the formal discrete measure of all zeros (or of a large window thereof),
\[
  \mu = \sum_{n\ge 1} \delta_{\gamma_n}.
\]
Then the ``ideal'' infinite-kernel resonance would involve
\[
  R_\infty(t) = \frac{1}{M}\sum_{n=1}^M K_\sigma(t-\gamma_n)
\]
for $M\to\infty$ in an appropriate sense, or, more conceptually, $(\mu * K_\sigma)(t)$.

Truncation introduces two types of spectral leakage:

\begin{enumerate}
  \item \textbf{Tail leakage.} Omitting zeros with $n>N$ eliminates their contribution, so the difference
  \[
    R_\infty(t) - R_N(t)
      = \frac{1}{M}\sum_{n>N} K_\sigma(t-\gamma_n) + \text{normalization error}
  \]
  represents missing tail energy. Even if the high zeros are individually small contributors, the cumulative effect may bias the kernel, especially for large $|t|$ away from the first $N$ zeros.

  \item \textbf{Bandwidth leakage.} Finite $\sigma$ smooths fine-scale spacing structure. Close zeros with spacing much smaller than $\sigma$ produce overlapping bands, effectively merging distinct spectral features into a single broad peak. This suppresses the high-frequency structure encoded in the pair correlation statistics and replaces it by a blurred approximation.
\end{enumerate}

From the point of view of Z-MOD, spectral leakage means that the auxiliary gradient derived from $R_N(t)$ encodes a low-rank, band-limited summary of zero statistics rather than the full fine-grained structure. As $N$ grows and $\sigma$ decreases, leakage can be reduced but at the cost of increased numerical stiffness and potential instability in the gradient field.

\subsection{Von Mangoldt Error-Term Regularization}

The classical explicit formula for the Chebyshev function $\psi(x)$ can be written schematically as
\[
  \psi(x) = x - \sum_{|\gamma_n| < T} \frac{x^{\rho_n}}{\rho_n} + E_T(x),
\]
where the sum runs over zeros with $|\gamma_n|<T$ and $E_T(x)$ collects the tail contributions and additional explicit terms.\footnote{Different normalizations and additional terms (such as contributions from the pole at $s=1$ and from trivial zeros) appear in standard references; here we isolate the truncated zero sum and its remainder.}

In a Z-MOD-inspired auxiliary loss, one might use the truncated oscillatory part as in $\mathcal{L}_{\mathrm{VM}}$ and introduce an error-term regularizer
\[
  \mathcal{R}_{\mathrm{err}}(\theta,\bm{\pi})
   = \lambda_{\mathrm{err}} \big|E_T(x(\theta,\bm{\pi}))\big|^2,
  \qquad x(\theta,\bm{\pi}) = e^{t(\theta,\bm{\pi})}.
\]
If explicit bounds of the form $|E_T(x)| \le B_T(x)$ are available on a relevant range of $x$, a practical surrogate is
\[
  \mathcal{R}_{\mathrm{err}}^{\mathrm{sur}}(\theta,\bm{\pi})
   = \lambda_{\mathrm{err}} B_T(x(\theta,\bm{\pi}))^2,
\]
which penalizes optimizer states where explicit-formula truncation is expected to be unreliable.

Differentiating the exact error-term penalty yields
\[
  \nabla_\theta \mathcal{R}_{\mathrm{err}}(\theta,\bm{\pi})
   = 2\lambda_{\mathrm{err}} E_T(x)\, E_T'(x)\, \nabla_\theta x(\theta,\bm{\pi}),
\]
with $\nabla_\theta x = x\,\nabla_\theta t$ and $\nabla_\theta t = a\nabla_\theta \tilde{\phi}(\theta)$. In the unclipped region,
\[
  \nabla_\theta x(\theta,\bm{\pi})
   = x(\theta,\bm{\pi})\, a\, \frac{2}{n}\theta,
\]
so
\[
  \nabla_\theta \mathcal{R}_{\mathrm{err}}(\theta,\bm{\pi})
   = 2\lambda_{\mathrm{err}} E_T(x(\theta,\bm{\pi}))\, E_T'(x(\theta,\bm{\pi}))\,
     x(\theta,\bm{\pi})\, a\, \frac{2}{n}\theta.
\]
Even if $E_T$ is not known exactly, using an upper bound $B_T$ or a numerically estimated surrogate for the magnitude of the omitted tail allows Z-MOD to damp regions where the truncated explicit formula is most questionable, stabilizing the influence of the von Mangoldt-based auxiliary terms.

\subsection{Finite-Width GUE Departures in Optimizer-Induced Spectra}

GUE universality is an asymptotic phenomenon: as the size of random matrices grows, their local eigenvalue statistics in the bulk converge to universal forms such as the sine kernel. In Z-MOD, several finite-width effects guarantee deviations from exact GUE behaviour:

\begin{itemize}
  \item Only finitely many zeros are used in the resonance layer.
  \item The sequence $\{t_k\}$ is not sampled from a stationary point process but generated by a non-stationary optimization dynamic that depends on the task loss.
  \item Gaussian smoothing with width $\sigma$ further blurs the spectrum, effectively convolving the discrete zero measure with a smooth kernel.
  \item The embedded heights $t_k$ are constrained by clipping of $\tilde{\phi}$ and by the optimizer's step-size schedule.
\end{itemize}

Thus any empirical pair correlation function $\widehat{R}_2(u)$ derived from normalized $\{t_k\}$ will, at best, approximate a smoothed and biased version of the GUE sine kernel on certain scales and windows. Rather than expecting exact universality, one should treat GUE statistics as a \emph{benchmark shape}: the closer $\widehat{R}_2(u)$ is to the sine kernel on appropriate windows, the more faithfully Z-MOD's induced spectral density reflects the structure of the underlying zero data.

From the standpoint of convergence, these finite-width departures are not problematic. The convergence of Z-MOD relies on the boundedness and Lipschitz properties of its kernelized gradients and on standard stochastic approximation conditions. GUE universality and Montgomery's PCC serve as design heuristics for the shape of the resonance kernel, not as analytic ingredients in the convergence proofs.

\begin{remark}
The correct division of labour is therefore:
\begin{itemize}
  \item Use analytic number theory---explicit formulas, zero data, and PCC/GUE intuition---to choose and tune the kernels $R_N(t)$ and explicit-formula-inspired terms in a way that is numerically stable and structurally interpretable.
  \item Use functional analysis and stochastic approximation theory to prove that, for the chosen kernels, the resulting gradient field is sufficiently regular (e.g., locally Lipschitz and bounded) to ensure convergence to stationary points under appropriate step-size and noise assumptions.
\end{itemize}
This separation ensures that any failures or refinements of RH or PCC affect the \emph{interpretation and shape} of the kernels but not the logical soundness of the convergence arguments.
\end{remark}

\section{Conclusion}

This document has formalized the Zeta Genius and Z-MOD developments into a coherent mathematical, algorithmic, and architectural specification. It introduced the multiplicity-aware state space, the canonical complex embedding, zeta-guided potentials, zero-data-informed resonance layers, multiplicity dynamics, PyTorch implementation sketches, and a phased ADR roadmap with benchmark protocols.

The combination of these elements constitutes a concrete prior art for zeta-guided, multiplicity-theoretic optimization frameworks. It is designed to support open scientific investigation of whether zeta-specific structure can genuinely affect optimization, while preventing later enclosure of the underlying ideas via patenting.

\appendix
\section*{Mathematical Appendix}

This appendix collects the explicit proofs, norm bounds, and auxiliary lemmas underlying the constructions used in the Z-MOD framework. The aim is to make the analytic foundations of the embedding, zeta potential, and zero-resonance layer completely explicit and self-contained.

Throughout, $n\in\mathbb{N}$ denotes the number of continuous parameters, $\mathcal{P}=\{p_1,\dots,p_m\}$ a fixed finite set of primes, and $\mathcal{M}=\{0,1\}^m$ the associated multiplicity space.

\section{Clipped Energy Functional and Embedding}

\subsection{Definitions}

Let $\theta\in\mathbb{R}^n$ and define
\[
  \phi(\theta) = \frac{1}{n}\|\theta\|^2 = \frac{1}{n}\sum_{i=1}^n \theta_i^2.
\]
Fix a clipping threshold $\Phi_{\max} > 0$ and define
\[
  \tilde{\phi}(\theta) = \min\big(\phi(\theta),\Phi_{\max}\big).
\]

For a multiplicity vector $\bm{\pi} = (\pi_{p_1},\dots,\pi_{p_m})\in\mathcal{M}$, define the canonical complex embedding
\begin{equation}
  s(\theta,\bm{\pi})
    = \sigma_0 + i\bigg(a\,\tilde{\phi}(\theta) + \sum_{p\in\mathcal{P}} \pi_p \log p\bigg),
  \label{A:eq:embedding}
\end{equation}
with fixed parameters $\sigma_0\in(0,1)$ and $a>0$.

\subsection{Basic properties}

\begin{lemma}[Continuity and boundedness of $\tilde{\phi}$]
For all $\theta\in\mathbb{R}^n$, $\tilde{\phi}(\theta)$ is continuous, non-negative, and satisfies
\[
  0 \le \tilde{\phi}(\theta) \le \Phi_{\max}.
\]
Moreover, $\tilde{\phi}$ is globally Lipschitz with constant $L_\phi = \frac{2}{\sqrt{n}}\Phi_{\max}^{1/2}$.
\end{lemma}

\begin{proof}
Continuity and non-negativity are immediate since $\phi(\theta)=\frac{1}{n}\|\theta\|^2$ is continuous and non-negative, and the minimum of continuous functions is continuous.

The bound $\tilde{\phi}(\theta)\le\Phi_{\max}$ follows directly from the definition $\tilde{\phi}(\theta) = \min(\phi(\theta),\Phi_{\max})$. Non-negativity is preserved under $\min(\cdot,\Phi_{\max})$.

To establish the Lipschitz bound, note that for any $\theta,\theta'\in\mathbb{R}^n$,
\[
  |\tilde{\phi}(\theta) - \tilde{\phi}(\theta')|
  \le |\phi(\theta) - \phi(\theta')|,
\]
since clipping cannot increase differences. Now
\[
  \phi(\theta) - \phi(\theta')
  = \frac{1}{n}(\|\theta\|^2 - \|\theta'\|^2)
  = \frac{1}{n}\langle \theta+\theta',\,\theta-\theta'\rangle,
\]
and hence
\[
  |\phi(\theta) - \phi(\theta')|
  \le \frac{1}{n}\|\theta+\theta'\|\cdot\|\theta-\theta'\|.
\]
Whenever $\phi(\theta)\le\Phi_{\max}$ and $\phi(\theta')\le\Phi_{\max}$, we have
\[
  \|\theta\|^2 \le n\Phi_{\max},\quad \|\theta'\|^2 \le n\Phi_{\max},
\]
whence $\|\theta+\theta'\|\le \|\theta\|+\|\theta'\| \le 2\sqrt{n\Phi_{\max}}$. Thus on the sublevel set $\{\theta:\phi(\theta)\le\Phi_{\max}\}$,
\[
  |\phi(\theta) - \phi(\theta')|
  \le \frac{2\sqrt{n\Phi_{\max}}}{n}\|\theta-\theta'\|
  = \frac{2}{\sqrt{n}}\Phi_{\max}^{1/2}\|\theta-\theta'\|.
\]
When one of $\phi(\theta)$ or $\phi(\theta')$ exceeds $\Phi_{\max}$, the difference is bounded by $\Phi_{\max}$ itself, while $\|\theta-\theta'\|\ge 0$; combining both cases yields a global Lipschitz constant $L_\phi = \frac{2}{\sqrt{n}}\Phi_{\max}^{1/2}$.
\end{proof}

\begin{lemma}[Lipschitz property of $s$]
For fixed $\bm{\pi}\in\mathcal{M}$, the map $\theta\mapsto s(\theta,\bm{\pi})$ is globally Lipschitz with respect to the Euclidean norm, with constant
\[
  L_s = a L_\phi = a\,\frac{2}{\sqrt{n}}\Phi_{\max}^{1/2}.
\]
Moreover, $|\mathrm{Im}\,s(\theta,\bm{\pi})|$ is bounded linearly in $\|\theta\|$ on sublevel sets of $\phi$.
\end{lemma}

\begin{proof}
The dependence of $s$ on $\theta$ is solely through $\tilde{\phi}(\theta)$ in its imaginary part; the real part $\sigma_0$ and the prime-shift sum depend only on $\bm{\pi}$. Hence
\[
  |s(\theta,\bm{\pi}) - s(\theta',\bm{\pi})|
  = a\,|\tilde{\phi}(\theta)-\tilde{\phi}(\theta')|
  \le a L_\phi \|\theta-\theta'\|,
\]
by the previous lemma.

On any set where $\phi(\theta)\le\Phi_{\max}$, we have $\tilde{\phi}(\theta)=\phi(\theta)$ and
\[
  |\mathrm{Im}\,s(\theta,\bm{\pi})|
  \le a\phi(\theta) + \sum_{p\in\mathcal{P}} |\log p|
  \le a\Phi_{\max} + \sum_{p\in\mathcal{P}} \log p,
\]
so the imaginary part is bounded on that set. For large $\|\theta\|$, $\tilde{\phi}$ is clipped at $\Phi_{\max}$ and the same bound holds.
\end{proof}

\section{Zeta Potential: Regularity and Bounds}

\subsection{Definition of the base zeta potential}

For $(\theta,\bm{\pi})\in\mathcal{S}$ define
\[
  s(\theta,\bm{\pi}) \text{ as in \eqref{A:eq:embedding}},
\]
and let $\varepsilon>0$ be fixed. The base zeta potential is
\begin{equation}
  \mathcal{Z}_0(\theta,\bm{\pi}) = \Re \log\big(|\zeta(s(\theta,\bm{\pi}))| + \varepsilon\big).
  \label{A:eq:zeta_potential}
\end{equation}

\subsection{Continuity and boundedness on compact sets}

\begin{lemma}[Continuity of $\mathcal{Z}_0$]
For each fixed $\bm{\pi}\in\mathcal{M}$, the map $\theta\mapsto \mathcal{Z}_0(\theta,\bm{\pi})$ is continuous on $\mathbb{R}^n$.
\end{lemma}

\begin{proof}
The composition of continuous maps is continuous. From the previous section, $\theta\mapsto s(\theta,\bm{\pi})$ is continuous. The Riemann zeta function $\zeta(s)$ is holomorphic on $\{s\in\mathbb{C}: \Re(s)\neq 1\}$, with a simple pole at $s=1$. By construction, $\sigma_0\in(0,1)$ is fixed and never equals 1, so $s(\theta,\bm{\pi})$ avoids the pole. Therefore, $\zeta(s(\theta,\bm{\pi}))$ is well-defined and continuous. The maps $z\mapsto |z|$, $x\mapsto x+\varepsilon$, and $x\mapsto \log(x)$ are continuous on $(0,\infty)$, hence the composition is continuous. Taking real part preserves continuity, so $\theta\mapsto \mathcal{Z}_0(\theta,\bm{\pi})$ is continuous.
\end{proof}

\begin{lemma}[Boundedness on compact sets]
Let $K\subset\mathbb{R}^n$ be compact. Then for fixed $\bm{\pi}\in\mathcal{M}$, there exists $M_K>0$ such that for all $\theta\in K$,
\[
  |\mathcal{Z}_0(\theta,\bm{\pi})| \le M_K.
\]
\end{lemma}

\begin{proof}
The image of $K$ under the continuous map $\theta\mapsto s(\theta,\bm{\pi})$ is a compact set $S_K\subset\{s : \Re(s)=\sigma_0\}$. Since $\zeta(s)$ is continuous on $S_K$ and $S_K$ avoids $s=1$, there exists $M>0$ such that $|\zeta(s)|\le M$ for all $s\in S_K$. Also $|\zeta(s)|\ge 0$. Thus for $\theta\in K$,
\[
  |\mathcal{Z}_0(\theta,\bm{\pi})|
  = \big|\Re\log(|\zeta(s(\theta,\bm{\pi}))|+\varepsilon)\big|
  \le \big|\log(|\zeta(s(\theta,\bm{\pi}))|+\varepsilon)\big|
  \le \max\{|\log \varepsilon|, |\log(M+\varepsilon)|\}.
\]
Set $M_K = \max\{|\log \varepsilon|, |\log(M+\varepsilon)|\}$.
\end{proof}

\subsection{Gradient and local Lipschitz bounds}

On any open region where $|\zeta(s)|>0$, i.e.\ $s$ avoids zeros of $\zeta$, we can differentiate $\mathcal{Z}_0$ with respect to $\phi$:

\[
  \frac{\partial \mathcal{Z}_0}{\partial \phi}
   = \Re\bigg(
      \frac{\zeta'(s)}{\zeta(s)}\cdot \frac{\partial s}{\partial \phi}
     \bigg)\cdot \frac{|\zeta(s)|}{|\zeta(s)|+\varepsilon}
   = \Re\bigg(
      \frac{\zeta'(s)}{\zeta(s)}\cdot i a
     \bigg)\cdot \frac{|\zeta(s)|}{|\zeta(s)|+\varepsilon}.
\]

Furthermore,
\[
  \nabla_\theta \phi(\theta) = \frac{2}{n}\theta,
\]
so for $\phi(\theta) < \Phi_{\max}$ we have
\begin{equation}
  \nabla_\theta \mathcal{Z}_0(\theta,\bm{\pi})
    = \Re\bigg(
        \frac{\zeta'(s(\theta,\bm{\pi}))}{\zeta(s(\theta,\bm{\pi}))}\cdot i a
      \bigg)\cdot \frac{|\zeta(s(\theta,\bm{\pi}))|}
                         {|\zeta(s(\theta,\bm{\pi}))|+\varepsilon}
       \cdot \frac{2}{n}\theta.
  \label{A:eq:grad_zeta_full}
\end{equation}
For $\phi(\theta)\ge\Phi_{\max}$, we define $\nabla_\theta\mathcal{Z}_0$ to be zero, reflecting the clipping in $\tilde{\phi}$.

\begin{lemma}[Local Lipschitzness of $\nabla_\theta \mathcal{Z}_0$]
Let $K\subset\mathbb{R}^n$ be a compact set such that for all $\theta\in K$ and $\bm{\pi}\in\mathcal{M}$, $s(\theta,\bm{\pi})$ avoids zeros of $\zeta$. Then there exists $L_K>0$ such that
\[
  \|\nabla_\theta \mathcal{Z}_0(\theta,\bm{\pi}) -
    \nabla_\theta \mathcal{Z}_0(\theta',\bm{\pi})\|
  \le L_K\|\theta-\theta'\|
\]
for all $\theta,\theta'\in K$.
\end{lemma}

\begin{proof}
On $K$, both $s(\theta,\bm{\pi})$ and $\theta$ are contained in compact sets, and $\zeta'/\zeta$ is continuous since $s$ avoids zeros. The maps
\[
  \theta \mapsto \theta, \quad
  \theta \mapsto \frac{\zeta'(s(\theta,\bm{\pi}))}{\zeta(s(\theta,\bm{\pi}))}, \quad
  \theta \mapsto \frac{|\zeta(s(\theta,\bm{\pi}))|}
                     {|\zeta(s(\theta,\bm{\pi}))|+\varepsilon}
\]
are all continuous on $K$ and attain finite maxima there. The expression for $\nabla_\theta\mathcal{Z}_0$ in \eqref{A:eq:grad_zeta_full} is a product of these bounded continuous functions and $\theta$, so it is Lipschitz on $K$ by standard results about smooth functions on compact domains and the equivalence of norms. Explicitly, one may bound $\|\nabla_\theta\mathcal{Z}_0(\theta,\bm{\pi})\|$ by $C_K\|\theta\|$ for some $C_K>0$, and differentiate again to see that the Jacobian is bounded, yielding the desired Lipschitz constant.
\end{proof}

\section{Zero-Resonance Layer: Regularity Proofs}

\subsection{Definition}

Let $\{\gamma_n\}_{n=1}^N$ be fixed real numbers (e.g.\ the first $N$ imaginary parts of non-trivial zeros of $\zeta$). For $t\in\mathbb{R}$ define
\[
  R(t) = \frac{1}{N}\sum_{n=1}^N \exp\bigg(-\frac{(t-\gamma_n)^2}{2\sigma^2}\bigg)
\]
for some $\sigma>0$.

Given $t(\theta,\bm{\pi}) = \mathrm{Im}\,s(\theta,\bm{\pi})$ and $\tilde{\phi}(\theta)$ as before, the zero-resonance auxiliary loss is
\[
  \mathcal{L}_{\mathrm{zero}}(\theta,\bm{\pi})
   = \alpha_{\mathrm{zero}}\;\tilde{\phi}(\theta)\; R(t(\theta,\bm{\pi})).
\]

\subsection{Derivative with respect to $\theta$}

We first compute $R'(t)$:
\[
  R'(t) = \frac{1}{N}\sum_{n=1}^N
    \exp\bigg(-\frac{(t-\gamma_n)^2}{2\sigma^2}\bigg)\cdot
    \bigg(-\frac{t-\gamma_n}{\sigma^2}\bigg).
\]

If we restrict to the region where $\phi(\theta)<\Phi_{\max}$, we have $\tilde{\phi}(\theta)=\phi(\theta)$ and
\[
  t(\theta,\bm{\pi}) = a\phi(\theta) + \sum_{p\in\mathcal{P}} \pi_p \log p.
\]
Then
\[
  \frac{\partial t}{\partial \phi} = a,\qquad
  \nabla_\theta \phi(\theta) = \frac{2}{n}\theta.
\]
By the chain rule,
\[
  \nabla_\theta t(\theta,\bm{\pi}) = \frac{\partial t}{\partial \phi}\nabla_\theta \phi(\theta)
    = a\cdot \frac{2}{n} \theta.
\]

Using the product rule,
\begin{align*}
  \nabla_\theta \mathcal{L}_{\mathrm{zero}}(\theta,\bm{\pi})
    &= \alpha_{\mathrm{zero}}
        \Big( \nabla_\theta \tilde{\phi}(\theta)\, R(t(\theta,\bm{\pi}))
             + \tilde{\phi}(\theta)\, R'(t(\theta,\bm{\pi}))\,
               \nabla_\theta t(\theta,\bm{\pi})\Big).
\end{align*}
Inside the unclipped region,
\[
  \nabla_\theta \tilde{\phi}(\theta) = \nabla_\theta \phi(\theta) = \frac{2}{n}\theta,
\]
and
\[
  \nabla_\theta t(\theta,\bm{\pi}) = a \frac{2}{n}\theta.
\]
Thus
\begin{equation}
  \nabla_\theta \mathcal{L}_{\mathrm{zero}}(\theta,\bm{\pi})
   = \alpha_{\mathrm{zero}} \frac{2}{n}\theta\, R(t(\theta,\bm{\pi}))
     + \alpha_{\mathrm{zero}} \tilde{\phi}(\theta) R'(t(\theta,\bm{\pi}))
        a \frac{2}{n}\theta.
  \label{A:eq:grad_Lzero}
\end{equation}
Outside the unclipped region ($\phi(\theta)\ge\Phi_{\max}$), one may set $\nabla_\theta\mathcal{L}_{\mathrm{zero}}=0$ by design.

\subsection{Lipschitzness on bounded sets}

\begin{lemma}[Boundedness of $R$ and $R'$]
For all $t\in\mathbb{R}$,
\[
  0 < R(t) \le 1, \qquad
  |R'(t)| \le \frac{C}{\sigma}
\]
for some constant $C>0$ depending only on $\{\gamma_n\}$ and $\sigma$.
\end{lemma}

\begin{proof}
Each term in the sum defining $R(t)$ is of the form $\exp(-\frac{(t-\gamma_n)^2}{2\sigma^2})$, which lies strictly between 0 and 1. Hence $0<R(t)\le 1$.

For $R'(t)$, note that for each $n$,
\[
  \bigg|\exp\bigg(-\frac{(t-\gamma_n)^2}{2\sigma^2}\bigg)\cdot
       \frac{t-\gamma_n}{\sigma^2}\bigg|
  \le \frac{1}{\sigma^2}\sup_{x\in\mathbb{R}} e^{-x^2/(2\sigma^2)}|x|
  = \frac{1}{\sigma^2} \cdot \sigma \sup_{y\in\mathbb{R}} e^{-y^2/2}|y|
  = \frac{K}{\sigma},
\]
for some $K>0$. Averaging over $n$ preserves the bound, so $|R'(t)|\le K/\sigma$.
\end{proof}

\begin{lemma}[Local Lipschitzness of $\nabla_\theta \mathcal{L}_{\mathrm{zero}}$]
Let $K\subset\mathbb{R}^n$ be a compact set on which $\phi(\theta)<\Phi_{\max}$ for all $\theta\in K$. Then there exists $L^{\mathrm{zero}}_K>0$ such that for all $\theta,\theta'\in K$ and all $\bm{\pi}\in\mathcal{M}$,
\[
  \|\nabla_\theta\mathcal{L}_{\mathrm{zero}}(\theta,\bm{\pi}) -
    \nabla_\theta\mathcal{L}_{\mathrm{zero}}(\theta',\bm{\pi})\|
  \le L^{\mathrm{zero}}_K \|\theta-\theta'\|.
\]
\end{lemma}

\begin{proof}
Within $K$, both $\tilde{\phi}(\theta)=\phi(\theta)$ and $t(\theta,\bm{\pi})$ are smooth functions of $\theta$, with derivatives bounded on $K$ due to compactness. Expression \eqref{A:eq:grad_Lzero} shows that $\nabla_\theta\mathcal{L}_{\mathrm{zero}}$ is a linear combination (with bounded coefficients) of $\theta$ and functions of $t(\theta,\bm{\pi})$. Since $R$ and $R'$ are bounded and smooth in $t$ and $t$ is Lipschitz in $\theta$, standard composition rules imply that $\nabla_\theta\mathcal{L}_{\mathrm{zero}}$ is Lipschitz on $K$. An explicit constant can be obtained by bounding the derivatives of $t$, $R$, and $R'$ on $K$, but the existence of $L^{\mathrm{zero}}_K$ suffices here.
\end{proof}

\section{Operator Norm Bounds for the Combined Gradient}

\subsection{Combined gradient}

The full objective (with zero-resonance but without explicit-formula cosine term) is
\[
  \mathcal{O}_{\mathrm{full}}(\theta,\bm{\pi})
   = f(\theta) + \lambda \mathcal{Z}_0(\theta,\bm{\pi})
     + \mathcal{L}_{\mathrm{zero}}(\theta,\bm{\pi})
     + \beta C(\bm{\pi}).
\]
The gradient with respect to $\theta$ in the unclipped region can be written as
\[
  \nabla_\theta\mathcal{O}_{\mathrm{full}}(\theta,\bm{\pi})
   = \nabla_\theta f(\theta)
     + \lambda \nabla_\theta \mathcal{Z}_0(\theta,\bm{\pi})
     + \nabla_\theta \mathcal{L}_{\mathrm{zero}}(\theta,\bm{\pi}).
\]

\subsection{Operator norm bounds on compact sets}

\begin{theorem}[Local Lipschitz bound for $\nabla_\theta \mathcal{O}_{\mathrm{full}}$]
Let $K\subset\mathbb{R}^n$ be a compact set such that:
\begin{enumerate}
  \item $\phi(\theta)<\Phi_{\max}$ for all $\theta\in K$;
  \item $s(\theta,\bm{\pi})$ avoids zeros of $\zeta$ for all $\theta\in K$ and $\bm{\pi}\in\mathcal{M}$;
  \item $\nabla_\theta f$ is Lipschitz on $K$ with constant $L^f_K$.
\end{enumerate}
Then there exists a constant $L^{\mathrm{full}}_K>0$ such that
\[
  \|\nabla_\theta\mathcal{O}_{\mathrm{full}}(\theta,\bm{\pi}) -
    \nabla_\theta\mathcal{O}_{\mathrm{full}}(\theta',\bm{\pi})\|
  \le L^{\mathrm{full}}_K \|\theta-\theta'\|
\]
for all $\theta,\theta'\in K$ and all $\bm{\pi}\in\mathcal{M}$.
\end{theorem}

\begin{proof}
From the assumptions, $\nabla_\theta f$ is Lipschitz on $K$ with constant $L^f_K$. From the previous lemmas, $\nabla_\theta\mathcal{Z}_0$ and $\nabla_\theta\mathcal{L}_{\mathrm{zero}}$ are Lipschitz on $K$ with constants $L^{\mathcal{Z}_0}_K$ and $L^{\mathrm{zero}}_K$ respectively (depending on $\lambda$, $\alpha_{\mathrm{zero}}$, and the choice of $N$, $\{\gamma_n\}$, $\sigma$). Thus
\begin{align*}
\|\nabla_\theta\mathcal{O}_{\mathrm{full}}(\theta,\bm{\pi}) -
  \nabla_\theta\mathcal{O}_{\mathrm{full}}(\theta',\bm{\pi})\|
&\le \|\nabla_\theta f(\theta) - \nabla_\theta f(\theta')\|
  + \lambda \|\nabla_\theta\mathcal{Z}_0(\theta,\bm{\pi}) -
              \nabla_\theta\mathcal{Z}_0(\theta',\bm{\pi})\|\\
&\quad + \|\nabla_\theta\mathcal{L}_{\mathrm{zero}}(\theta,\bm{\pi}) -
          \nabla_\theta\mathcal{L}_{\mathrm{zero}}(\theta',\bm{\pi})\| \\
&\le \big(L^f_K + \lambda L^{\mathcal{Z}_0}_K + L^{\mathrm{zero}}_K\big)\|\theta-\theta'\|.
\end{align*}
Setting $L^{\mathrm{full}}_K = L^f_K + \lambda L^{\mathcal{Z}_0}_K + L^{\mathrm{zero}}_K$ yields the claim.
\end{proof}

\section{Remarks on Numerical Stability and Clipping}

The theoretical results above justify the use of clipping in $\tilde{\phi}(\theta)$ and in $\mathrm{Im}\,s$ from the standpoint of preserving Lipschitzness and boundedness of the zeta potentials and their gradients on regions of interest.

From a practical perspective:
\begin{itemize}[leftmargin=*]
  \item Clipping ensures $\mathrm{Im}\,s$ remains within a manageable numerical range for evaluating $\zeta(s)$ and $R(t)$.
  \item The auxiliary loss $\mathcal{L}_{\mathrm{zero}}$ inherits good regularity properties from the Gaussian kernel structure of $R$.
  \item Operator norm bounds on the Hessians (or Jacobians of the gradients) on compact sets can be obtained by differentiating the expressions above, though explicit closed forms are increasingly cumbersome.
\end{itemize}

Nevertheless, the key property required for the optimizer design is that $\nabla_\theta\mathcal{O}_{\mathrm{full}}$ is locally Lipschitz in $\theta$ on relevant parameter regions, which the lemmas and theorem provide.

\nocite{*}
\bibliographystyle{unsrt}  
\bibliography{references}  


\end{document}
